% ======================================================= %
% Document: TEMPLATE FOR RESPONSES TO REVIEWERS
% Author: Andrea Ballatore
% Date: Jan 7, 2013
% Source: https://raw.githubusercontent.com/ucd-spatial/Datasets/master/tex_response_to_reviewers_template/responses_to_reviewers.tex
% Modified by Eduard Szöcs, 10.03.2015
% ======================================================= %
\documentclass[12pt]{article}

% packages
\usepackage{xr-hyper}
\usepackage{hyperref}
\externaldocument[ms-]{ant_cactus_SUBMISSION2}

\usepackage{graphicx}
\usepackage{url}
\usepackage[usenames,dvipsnames]{xcolor}
\usepackage{color}
\definecolor{mygray}{gray}{0.6}
\usepackage[utf8]{inputenc}
\usepackage[onehalfspacing]{setspace}
\usepackage[
round,	%(defaultage in the main file and \input ) for round parentheses;
colon,	% (default) to separate multiple citations with colons;
authoryear,% (default) for author-year citations;
sort,		% orders multiple citations into the sequence in which they
]{natbib}
\usepackage[%disable
]{todonotes}

\usepackage{anysize}
\marginsize{2.5cm}{2.5cm}{1.5cm}{2.5cm}

% macros
% add a counter
\newcounter{cN}
\setcounter{cN}{0}

\newcommand{\comment}[1]{
	\vspace{2em}
	\refstepcounter{cN} % incrment counter
	\noindent \hangindent=0em \textbf{\textcolor{Maroon}{\uline{Comment \thecN}:~}}\emph{``#1''}
}

\newcommand{\response}[1]{
	\\[0.25em]
	\hangindent=2.3em \textbf{\textcolor{NavyBlue}{\uline{Response}:~}}#1
}

\newcommand{\revise}[1]{{\color{blue}{#1}}}
\newcommand{\tom}[2]{{\color{red}{#1}}\footnote{\textit{\color{red}{#2}}}} 
\usepackage[normalem]{ulem}
\definecolor{darkred}{rgb}{1,.6,.6}
\DeclareRobustCommand\problemline{\bgroup\markoverwith{\textcolor{darkred}{\rule[-0.9ex]{4pt}{3pt}}}\ULon}
\DeclareRobustCommand{\problem}[1]{\problemline{#1}} % soul
\setcounter{secnumdepth}{-1}


\begin{document}
	
	\title{Manuscript 62524R1 --- Response to Reviewers}
	\date{}
	\maketitle
	
	\noindent To the Editorial Board of \emph{The American Naturalist},
	
	Thank you for the``contingent accept'' decision regarding our manuscript, ``Demographic consequences of partner diversity and turnover in a multi-species ant-plant mutualism.'' We are grateful to Dr. Frederickson and the reviewers for their continued support and constructive feedback.
	
	In this final revision, we have addressed the remaining minor textual and figure concerns raised by Reviewer 3. Furthermore, we have performed a comprehensive overhaul of our data repository and analysis scripts to meet the high standards of reproducibility required by the journal, as outlined by the Data Editor. 
	
	Specifically, we have:
	\begin{enumerate}
		\item Consolidated all code and data into a single, streamlined repository.
		\item Replaced absolute file paths with relative paths to ensure the code runs on any machine.
		\item Improved the README file with detailed instructions, version requirements, and a roadmap for reproducing all figures and tables.
		\item Refined Figure 5 and other graphical elements for better resolution and clarity.
	\end{enumerate}
	
	Detailed responses to the specific comments follow below.
	
	\vspace{2em}
	\hfill Sincerely,
	\hfill The Authors
	
	\newpage
	
	% ======================================================= %
	\section{Response to Reviewer 3}
	
	\comment{L460: A central question posed by the authors is how vital rates fluctuate across years... It could be helpful for later understanding portfolio effects if the authors explicitly compare the variance of year-specific ant effects on the host in the text of this section...}
	\response{This is a fascinating question! Unfortunately we are not well-equipped to answer it with our current analyses. For statistical efficiency, our models assume that all ant states share the same inter-annual variance (see Equations 4-7 and corresponding text). While we allow year effects to be ant-specific (which facilitated the portfolio effect analysis), the random deviates are assumed to come from the same distribution of inter-annual variability. The alternative -- ant-specific temporal variances -- would be very interesting to explore but also very data-hungry, pushing the limits of what we can learn from even this unusually large data set. Hence our decision to make a simplifying assumption that improves model efficiency by ``borrowing information'' across ant states to estimate an aggregate temporal variance. 
	\\
	\\
	To clarify this issue for other readers, we now make this assumption more explicit in our presentation of the statistical models (see line \ref{ms-res:ant_year_correlations}).}
	
	\comment{I was very interested in the scenario in which dominant ant species did not competitively exclude the rest under which a sampling effect emerged... I think that it would improve the manuscript to briefly reference this scenario and the intuition behind why a sampling effect emerged...}
	\response{We agree that this helps clarify the ecological conditions necessary for a sampling effect. We have added a discussion of this scenario and the underlying intuition in the Results (lines \ref{ms-res:sampling_effect}) and added a brief bridging sentence in the Discussion (line \ref{ms-sampling_effect2}) to highlight why the sampling effect was absent in our main model but present in this alternative scenario.}
	
	\comment{L95 and L97: Suggestions regarding 'opportunity cost' and 'relative fitness'.}
	\response{We have adopted both suggestions. We changed ``missed opportunity cost'' to ``opportunity cost'' throughout the manuscript (lines \ref{ms-res:L95-1} and \ref{ms-res:L95-2}) and specified ``relative fitness'' where appropriate to ensure technical accuracy regarding fitness benefits (line \ref{ms-res:lL5-1}).}
	
	\comment{L297: It would be helpful to add underbrackets labeling the major components of equation 3.}
	\response{We agree that this would improve the clarity of the equation for the readers and have added in underbrackets with descriptive labels for this equation (eqn 3).}
	
	\comment{L413, 431 \& ,541: references to specific appendix sections remain missing throughout the manuscript.}
	\response{We have corrected these  so the references work properly with the appendix.}
	
	\comment{L479: Figures 1, 2, and 5 have low resolution.}
	\response{We have resolved these figure resolution issues (lines \ref{ms-fig:resolution1}, \ref{ms-fig:resolution2}, and \ref{ms-fig:resolution3}).}
	
	
	\comment{Figure 4a: the color legend does not match the palette of the colors in the correlation plot.}
	\response{We have corrected this error (line \ref{ms-fig:corrplot_colors}).}
	
	\comment{Figure 5: the margins of the figure's panels overlap, causing part of the panel label to be lost for each panel. Every panel of this figure is labeled with a b.}
	\response{We have corrected the margin overlaps (line \ref{ms-fig:resolution3}).}
	
	% ======================================================= %
\section{Response to the Data Editor}

\comment{Several scripts include absolute file paths (e.g., setwd("/Users/alicampbell/Documents/GitHub/ant\_cactus\_demography") which will not work on other systems.}
\response{We have removed all `setwd()` calls and replaced absolute pathways with relative pathways. This should fix any issues with absolute pathways on other machines.}

\comment{Many scripts reference `.Rds` files that were not included in the materials, making reproduction difficult or impossible.}
\response{We have also eliminated the need for .Rds files and instead suggest that the model outputs previously stored in these files be generated through the code in file 02\_cholla\_ant\_IPM\_vital\_rates.R}

\comment{The README is currently insufficient for guiding users through the workflow. It lacks detailed instructions, file descriptions, and information about software versions and script dependencies.}
\response{The README has been extensively updated. It now includes much more detail about the workflow, files, software versions, and script dependencies.}

\comment{**Consolidate all materials** into a single, self-contained archive or repository.}
\response{We have done this. The respository has been reorganized and now includes all necessary materials for running this code. }

\comment{**Test your code** on a clean R session or a different machine to ensure it runs as expected. This will help ensure that all no analyses rely on the paths set by the code.}
\response{The code was tested on a clean R session to ensure the workflow is as expected and nothing is unloaded or sourced from a different location.}
	
\end{document}